\documentclass[12pt]{article}
\usepackage[T1]{fontenc}
\usepackage[utf8]{inputenc}
\usepackage{lmodern}
\usepackage{textcomp}
\usepackage{enumitem}
\usepackage{geometry}
\usepackage{graphicx}
\usepackage{amsmath, amssymb}
\geometry{margin=1in}
\begin{document}
\begin{center}
Biology - Graduate Level Assessment 11 \
Total Marks: 40 \
Answer all questions.
\end{center}

\section*{Multiple Choice (10 marks)}
\begin{enumerate}[label=\arabic*.]
\item Two xylem plant cell types that provide support and conduct water and minerals are the
\begin{enumerate}[label=(\alph*)]
    \item parenchyma and companion cells
    \item tracheids and vessel elements
    \item parenchyma and collenchyma
    \item sieve tube members and companion cells
    \item collenchyma and vessel elements
\end{enumerate}
\item Hemoglobin is a molecule that binds to both O 2 and CO 2. There is an allosteric relationship between the concentrations of O 2 and CO 2. Hemoglobin's affinity for O 2
\begin{enumerate}[label=(\alph*)]
    \item increases as H+ concentration increases
    \item decreases as blood pH decreases
    \item decreases in resting muscle tissue
    \item increases as blood pH decreases
    \item increases as O 2 concentration decreases
\end{enumerate}
\item Explain the mechanism of the genetic determination of sex in man .
\begin{enumerate}[label=(\alph*)]
    \item The temperature during embryonic development determines the sex.
    \item The presence of an X or Y chromosome in the mother determines the sex of the offspring.
    \item The presence of the X chromosome determines the sex.
    \item The sex is determined by environmental factors after birth.
    \item The presence of two Y chromosomes determines that an individual will be male.
\end{enumerate}
\item What unusual feature of the cycads makes them are remarkable group among the seed plants?
\begin{enumerate}[label=(\alph*)]
    \item They have flowers that bloom for several years before producing seeds
    \item They have a unique leaf structure
    \item They can survive in extremely cold climates
    \item They produce large seeds
    \item Their seeds can remain dormant for up to a century before germinating
\end{enumerate}
\item The energy given up by electrons as they move through the electron transport chain is used to
\begin{enumerate}[label=(\alph*)]
    \item produce ATP
    \item create carbon dioxide
    \item make FADH 2
    \item synthesize proteins
    \item make glucose
\end{enumerate}
\end{enumerate}

\section*{True / False (10 marks)}
\textit{Answer True or False}
\begin{enumerate}[label=\arabic*.]
\setcounter{enumi}{5}
\item True or False: Is the advanced age a contraindication to GERD laparoscopic surgery?
\item True or False: Should ascitis volume and anthropometric measurements be estimated in hospitalized alcoholic cirrotics?
\item True or False: Perioperative care in an animal model for training in abdominal surgery: is it necessary a preoperative fasting?
\item True or False: Does Sensation Return to the Nasal Tip After Microfat Grafting?
\item True or False: Is fetal anatomic assessment on follow-up antepartum sonograms clinically useful?
\end{enumerate}

\section*{Long Form Response (20 marks)}
\textit{Answer each question with a clear and complete explanation.}
\begin{enumerate}[label=\arabic*.]
\setcounter{enumi}{10}
\item Discuss the two fundamental types of cells identified by biologists, highlighting their key characteristics, differences, and significance in biological systems.
\item Discuss the essential characteristics that define living organisms and explain how these attributes differentiate living matter from non-living chemical components.
\end{enumerate}

\vfill
\noindent\textit{End of Paper}
\end{document}